\section{Module \texttt{toggle.c}}

Le module \texttt{toggle.c} couvre les composants binaires du wrapper : case à cocher et ligne composée avec interrupteur. Ces deux widgets reposent sur des contrôles GTK simples, mais le wrapper leur ajoute une configuration déclarative homogène.

\subsection{Fonction \texttt{create\_checkbox}}

\begin{lstlisting}[language=C]
GtkWidget *create_checkbox(const checkbox_config *config) {
    GtkWidget *checkbox = gtk_check_button_new_with_label(config && config->label ? config->label : "");

    if (config == NULL) {
        return checkbox;
    }

    gtk_check_button_set_active(GTK_CHECK_BUTTON(checkbox), config->active);
    gtk_check_button_set_inconsistent(GTK_CHECK_BUTTON(checkbox), config->inconsistent);
    apply_widget_style(checkbox, &config->style);

    if (config->on_toggled != NULL) {
        g_signal_connect(checkbox, "toggled", G_CALLBACK(config->on_toggled), config->user_data);
    }

    return checkbox;
}
\end{lstlisting}

\subsection{Fonction \texttt{create\_switch\_row}}

Cette seconde fonction construit un composant composite : un label à gauche, un \texttt{GtkSwitch} à droite, puis un éventuel callback sur la propriété \texttt{active}.

\begin{lstlisting}[language=C]
GtkWidget *create_switch_row(const switch_config *config, GtkWidget **out_switch) {
    GtkWidget *row = gtk_box_new(GTK_ORIENTATION_HORIZONTAL, 10);
    GtkWidget *label = gtk_label_new(config && config->label ? config->label : "");
    GtkWidget *sw = gtk_switch_new();

    gtk_widget_set_hexpand(label, TRUE);
    gtk_widget_set_halign(label, GTK_ALIGN_START);
    gtk_widget_set_halign(sw, GTK_ALIGN_END);
    gtk_box_append(GTK_BOX(row), label);
    gtk_box_append(GTK_BOX(row), sw);

    if (config != NULL) {
        gtk_switch_set_active(GTK_SWITCH(sw), config->active);
        gtk_switch_set_state(GTK_SWITCH(sw), config->state);
        apply_widget_style(row, &config->style);

        if (config->on_active_changed != NULL) {
            g_signal_connect(sw, "notify::active", G_CALLBACK(config->on_active_changed), config->user_data);
        }
    }

    if (out_switch != NULL) {
        *out_switch = sw;
    }

    return row;
}
\end{lstlisting}

\subsection{APIs GTK utilisées}

\begin{itemize}
    \item \textbf{\texttt{gtk\_check\_button\_new\_with\_label()}} : Initialise une case à cocher associée à un texte descriptif. Le module l'utilise pour créer des contrôles de sélection binaire simples et lisibles.
    \item \textbf{\texttt{gtk\_check\_button\_set\_inconsistent()}} : Permet de placer visuellement la case à cocher dans un état tiers, indiquant une sélection partielle ou indéterminée. Le wrapper expose cette fonctionnalité pour gérer des logiques de sélection complexes.
    \item \textbf{\texttt{gtk\_switch\_new()}} : Instancie un widget d'interrupteur coulissant, typique des interfaces modernes. Elle est au cœur de la fonction \texttt{create\_switch\_row} pour proposer une interaction intuitive sur les écrans tactiles ou avec la souris.
    \item \textbf{\texttt{gtk\_switch\_set\_active()}} : Modifie l'état booléen du switch sans déclencher de signaux de changement d'état. Le wrapper l'utilise pour appliquer l'état initial défini par le développeur lors de la création du widget.
    \item \textbf{\texttt{gtk\_box\_append()}} : Insère un widget à la suite des autres dans un conteneur \texttt{GtkBox}. Dans le module, elle sert à construire dynamiquement la ligne composite regroupant le libellé et l'interrupteur.
    \item \textbf{\texttt{g\_signal\_connect()}} : Raccorde les signaux de changement d'état (comme \texttt{notify::active} pour le switch) aux callbacks applicatifs. Elle assure la réactivité de l'application aux changements d'état opérés par l'utilisateur sur les contrôles de bascule.
\end{itemize}

\newpage
